%% ============================================================================
%% Section: The Causal Query Language
%% Included from main.tex via %% ============================================================================
%% Section: The Causal Query Language
%% Included from main.tex via %% ============================================================================
%% Section: The Causal Query Language
%% Included from main.tex via %% ============================================================================
%% Section: The Causal Query Language
%% Included from main.tex via \input{sections/language}
%% Requires the `causaldsl' listings style and the \TODO/\NOTE macros from main.tex.
%% Table label tab:operators is referenced from sections/background.tex (\S2.2).
%% Operator->procedure mapping is faithful to dsl/.../exec/CypherTranslator.java
%% and plugin/.../api/ProcedureNames.java.
%% ============================================================================
\section{The Causal Query Language}\label{sec:lang}

Requirement R1 asks for a declarative interface in which each catalog query is a
single statement, phrased close to how an analyst states a causal question, that
compiles to physical operators over the graph store. This section presents that
language: its surface (\S\ref{sec:lang:surface}), the mapping from each operator
to the physical operator that answers it (\S\ref{sec:lang:compile},
Table~\ref{tab:operators}), and the model registry that lets a single database
hold several causal worlds at once (\S\ref{sec:lang:registry}). The physical
operators themselves---and why they are fast---are the subject of
\S\ref{sec:index}; here we are concerned only with the interface and its
compilation.

\subsection{Design and surface syntax}\label{sec:lang:surface}

The language is a small statement-oriented DSL. A program is a sequence of
statements; each statement is one of three kinds. \emph{Definition} statements
declare variables and edges and thereby write the causal model into the graph.
\emph{Structural} statements ask the catalog queries of \S\ref{sec:bg:catalog}:
reachability, path roles, and identification. \emph{Quantitative} statements
express interventions and counterfactuals over a parameterized model
(\S\ref{sec:quant}). Two design choices follow directly from R1. First, an
operator names a causal concept, not a traversal: an analyst writes
\lstinline[style=causaldsl]|CONFOUNDERS OF T AND Y| rather than a path pattern,
so the statement reads as the question and its compilation to a graph operation
is the system's concern, not the user's. Second, every statement is
self-contained and refers to variables by their canonical
\textsf{id}, so a query is portable across any database holding a model with
those variables.

Definition statements come in a standalone and a block form. A bare
\lstinline[style=causaldsl]|DEFINE CAUSAL VARIABLE| introduces a single variable
with its observation status; a \lstinline[style=causaldsl]|DEFINE CAUSAL DAG|
block declares a model as a set of edges, optionally carrying a structural
coefficient per edge for the quantitative layer:

\begin{lstlisting}[style=causaldsl]
DEFINE CAUSAL DAG CovidModel {
  covid19_burden    -> interventions
  covid19_burden    -> exposure_sarscov2
  interventions     -> mobility
  mobility          -> exposure_sarscov2
  exposure_sarscov2 -> onset_covid19
  onset_covid19     -> reported_cases
}
\end{lstlisting}

\noindent
Against the model so defined, the structural questions of our running example
(\S\ref{sec:bg:example}) are each one statement:

\begin{lstlisting}[style=causaldsl]
CONFOUNDERS OF interventions AND reported_cases
BACKDOOR PATHS BETWEEN interventions AND reported_cases
IDENTIFY EFFECT OF interventions ON reported_cases
  RETURN ADJUSTMENT SETS EXPLAIN
\end{lstlisting}

\noindent
The first returns \textsf{covid19\_burden}; the second reports the backdoor path
it opens; the third decides identifiability and, with
\lstinline[style=causaldsl]|RETURN ADJUSTMENT SETS|, returns the valid adjustment
sets, the optional \lstinline[style=causaldsl]|EXPLAIN| attaching a
human-readable justification. The reachability and path-role operators
(\lstinline[style=causaldsl]|CAUSES OF|, \lstinline[style=causaldsl]|EFFECTS OF|,
\lstinline[style=causaldsl]|ANCESTORS OF|, \lstinline[style=causaldsl]|DESCENDANTS OF|,
\lstinline[style=causaldsl]|MEDIATORS BETWEEN|, \lstinline[style=causaldsl]|COLLIDERS BETWEEN|)
follow the same one-statement-per-question shape. The full grammar is given in
ANTLR form in the artifact; Table~\ref{tab:operators} lists the operators that
matter for what follows.

\subsection{Compilation to physical operators}\label{sec:lang:compile}

A statement is parsed to an AST and then compiled, statement by statement, to one
or more Cypher units that execute inside the DBMS. The compiler emits exactly two
kinds of unit. A \emph{definition} statement compiles to Cypher \emph{writes}
that maintain the schema of \S\ref{sec:bg:dag}: a \lstinline[style=causaldsl]|DEFINE|
becomes a \texttt{MERGE} of the \textsf{Variable} nodes and \textsf{CAUSES} edges
it introduces (a single batched \texttt{UNWIND} for the edge set), followed by a
call to \texttt{causal.compile} that refreshes the compiled snapshot
(\S\ref{sec:compile}). A \emph{structural} statement compiles to a read: a
\texttt{CALL} of the \texttt{causal.*} procedure that answers it, which consults
the snapshot rather than re-traversing the stored graph. This two-target
compilation is what keeps the model \emph{in} the database: the language never
holds the graph in client memory, and the same authoritative snapshot serves
every reader.

\begin{table}[t]
  \centering
  \caption{Language operators, their graphical semantics, and the physical
  operator each compiles to. Reachability and path-role operators map to
  \texttt{causal.*} procedures over the compiled snapshot (\S\ref{sec:index});
  d-separation is exposed as \texttt{causal.dsep} and is the inner test that
  identification issues repeatedly. Definition statements compile to Cypher
  writes plus a snapshot recompile.}
  \label{tab:operators}
  \small
  \setlength{\tabcolsep}{4pt}
  \begin{tabular}{@{}p{0.30\linewidth} p{0.34\linewidth} p{0.28\linewidth}@{}}
    \toprule
    \textbf{DSL statement} & \textbf{Graphical meaning} & \textbf{Physical operator} \\
    \midrule
    \multicolumn{3}{@{}l}{\textit{Definition}}\\
    \lstinline[style=causaldsl]|DEFINE ... DAG| & write / update model & \texttt{MERGE} + \texttt{causal.compile} \\
    \addlinespace[2pt]
    \multicolumn{3}{@{}l}{\textit{Reachability}}\\
    \lstinline[style=causaldsl]|CAUSES OF X|       & $\mathrm{pa}(X)$ & \texttt{causal.parents} \\
    \lstinline[style=causaldsl]|EFFECTS OF X|      & $\mathrm{ch}(X)$ & \texttt{causal.children} \\
    \lstinline[style=causaldsl]|ANCESTORS OF X|    & $\mathrm{an}(X)$ & \texttt{causal.ancestors} \\
    \lstinline[style=causaldsl]|DESCENDANTS OF X|  & $\mathrm{de}(X)$ & \texttt{causal.descendants} \\
    \addlinespace[2pt]
    \multicolumn{3}{@{}l}{\textit{Path roles}}\\
    \lstinline[style=causaldsl]|MEDIATORS ...|   & on a directed $T\!\to\!Y$ path & \texttt{causal.mediators} \\
    \lstinline[style=causaldsl]|COLLIDERS ...|   & $\to Z \leftarrow$ on a path & \texttt{causal.colliders} \\
    \lstinline[style=causaldsl]|CONFOUNDERS ...| & $\mathrm{an}(T)\cap\mathrm{an}(Y)$ & \texttt{causal.confounders} \\
    \lstinline[style=causaldsl]|BACKDOOR PATHS ...| & path with edge into $T$ & \texttt{causal.backdoorVariables} \\
    \addlinespace[2pt]
    \multicolumn{3}{@{}l}{\textit{Independence \& identification}}\\
    (inner test) & $X\!\perp\!Y\mid Z$ (d-sep) & \texttt{causal.dsep} \\
    \lstinline[style=causaldsl]|IDENTIFY EFFECT ...| & backdoor criterion, adj.\ sets & \texttt{causal.adjustmentCandidate} \\
    \bottomrule
  \end{tabular}
\end{table}

Structural operators map to procedures one-to-one. The reachability operators
compile to \texttt{causal.parents}, \texttt{causal.children},
\texttt{causal.ancestors}, and \texttt{causal.descendants}; the path-role
operators to \texttt{causal.mediators}, \texttt{causal.colliders}, and
\texttt{causal.confounders}; \lstinline[style=causaldsl]|BACKDOOR PATHS| to
\texttt{causal.backdoorVariables}; and \lstinline[style=causaldsl]|IDENTIFY EFFECT|
to \texttt{causal.adjustmentCandidate}, which returns whether an adjustment set
exists, the candidate sets, and---under \lstinline[style=causaldsl]|EXPLAIN|---an
explanation. D-separation is not a top-level keyword but the inner test that
identification issues repeatedly; it is exposed as the \texttt{causal.dsep}
procedure so that the cost of the traversal-intensive query (R2) is borne by a
single indexed operator (\S\ref{sec:index}) rather than by client-side path
enumeration. The quantitative statements (\lstinline[style=causaldsl]|INTERVENE|,
\lstinline[style=causaldsl]|COUNTERFACTUAL|) are handled by the quantitative
layer of \S\ref{sec:quant}; numeric intervention compiles to
\texttt{causal.linearGaussianInference}, while world construction is resolved in
the registry described next.

\subsection{Model registry and causal worlds}\label{sec:lang:registry}

A causal database is rarely a single fixed graph: a curator maintains several
models, and an intervention $do(X{=}x)$ defines a \emph{derived} model---the same
DAG with the edges into $X$ removed---that an analyst wants to name, query, and
compare against the original. The language therefore addresses models by name.
\lstinline[style=causaldsl]|DEFINE MODEL| registers a base model;
\lstinline[style=causaldsl]|USE MODEL| makes a named model (base or derived) the
active graph against which subsequent structural statements run; and
\lstinline[style=causaldsl]|RESET| returns to the base model while leaving derived
worlds registered and addressable. Context switching is resolved entirely in the
registry and emits no Cypher, so moving between worlds costs a name lookup rather
than a re-read of the graph. This registry is also the seam at which the
quantitative layer attaches: an \lstinline[style=causaldsl]|INTERVENE| statement
materializes its derived world once and registers it under a canonical name, so
that a later \lstinline[style=causaldsl]|USE MODEL| over that world reuses the
same compiled snapshot. The mechanics of snapshotting a model---base or
derived---into the layout the procedures read is the subject of the next section.

%% Requires the `causaldsl' listings style and the \TODO/\NOTE macros from main.tex.
%% Table label tab:operators is referenced from sections/background.tex (\S2.2).
%% Operator->procedure mapping is faithful to dsl/.../exec/CypherTranslator.java
%% and plugin/.../api/ProcedureNames.java.
%% ============================================================================
\section{The Causal Query Language}\label{sec:lang}

Requirement R1 asks for a declarative interface in which each catalog query is a
single statement, phrased close to how an analyst states a causal question, that
compiles to physical operators over the graph store. This section presents that
language: its surface (\S\ref{sec:lang:surface}), the mapping from each operator
to the physical operator that answers it (\S\ref{sec:lang:compile},
Table~\ref{tab:operators}), and the model registry that lets a single database
hold several causal worlds at once (\S\ref{sec:lang:registry}). The physical
operators themselves---and why they are fast---are the subject of
\S\ref{sec:index}; here we are concerned only with the interface and its
compilation.

\subsection{Design and surface syntax}\label{sec:lang:surface}

The language is a small statement-oriented DSL. A program is a sequence of
statements; each statement is one of three kinds. \emph{Definition} statements
declare variables and edges and thereby write the causal model into the graph.
\emph{Structural} statements ask the catalog queries of \S\ref{sec:bg:catalog}:
reachability, path roles, and identification. \emph{Quantitative} statements
express interventions and counterfactuals over a parameterized model
(\S\ref{sec:quant}). Two design choices follow directly from R1. First, an
operator names a causal concept, not a traversal: an analyst writes
\lstinline[style=causaldsl]|CONFOUNDERS OF T AND Y| rather than a path pattern,
so the statement reads as the question and its compilation to a graph operation
is the system's concern, not the user's. Second, every statement is
self-contained and refers to variables by their canonical
\textsf{id}, so a query is portable across any database holding a model with
those variables.

Definition statements come in a standalone and a block form. A bare
\lstinline[style=causaldsl]|DEFINE CAUSAL VARIABLE| introduces a single variable
with its observation status; a \lstinline[style=causaldsl]|DEFINE CAUSAL DAG|
block declares a model as a set of edges, optionally carrying a structural
coefficient per edge for the quantitative layer:

\begin{lstlisting}[style=causaldsl]
DEFINE CAUSAL DAG CovidModel {
  covid19_burden    -> interventions
  covid19_burden    -> exposure_sarscov2
  interventions     -> mobility
  mobility          -> exposure_sarscov2
  exposure_sarscov2 -> onset_covid19
  onset_covid19     -> reported_cases
}
\end{lstlisting}

\noindent
Against the model so defined, the structural questions of our running example
(\S\ref{sec:bg:example}) are each one statement:

\begin{lstlisting}[style=causaldsl]
CONFOUNDERS OF interventions AND reported_cases
BACKDOOR PATHS BETWEEN interventions AND reported_cases
IDENTIFY EFFECT OF interventions ON reported_cases
  RETURN ADJUSTMENT SETS EXPLAIN
\end{lstlisting}

\noindent
The first returns \textsf{covid19\_burden}; the second reports the backdoor path
it opens; the third decides identifiability and, with
\lstinline[style=causaldsl]|RETURN ADJUSTMENT SETS|, returns the valid adjustment
sets, the optional \lstinline[style=causaldsl]|EXPLAIN| attaching a
human-readable justification. The reachability and path-role operators
(\lstinline[style=causaldsl]|CAUSES OF|, \lstinline[style=causaldsl]|EFFECTS OF|,
\lstinline[style=causaldsl]|ANCESTORS OF|, \lstinline[style=causaldsl]|DESCENDANTS OF|,
\lstinline[style=causaldsl]|MEDIATORS BETWEEN|, \lstinline[style=causaldsl]|COLLIDERS BETWEEN|)
follow the same one-statement-per-question shape. The full grammar is given in
ANTLR form in the artifact; Table~\ref{tab:operators} lists the operators that
matter for what follows.

\subsection{Compilation to physical operators}\label{sec:lang:compile}

A statement is parsed to an AST and then compiled, statement by statement, to one
or more Cypher units that execute inside the DBMS. The compiler emits exactly two
kinds of unit. A \emph{definition} statement compiles to Cypher \emph{writes}
that maintain the schema of \S\ref{sec:bg:dag}: a \lstinline[style=causaldsl]|DEFINE|
becomes a \texttt{MERGE} of the \textsf{Variable} nodes and \textsf{CAUSES} edges
it introduces (a single batched \texttt{UNWIND} for the edge set), followed by a
call to \texttt{causal.compile} that refreshes the compiled snapshot
(\S\ref{sec:compile}). A \emph{structural} statement compiles to a read: a
\texttt{CALL} of the \texttt{causal.*} procedure that answers it, which consults
the snapshot rather than re-traversing the stored graph. This two-target
compilation is what keeps the model \emph{in} the database: the language never
holds the graph in client memory, and the same authoritative snapshot serves
every reader.

\begin{table}[t]
  \centering
  \caption{Language operators, their graphical semantics, and the physical
  operator each compiles to. Reachability and path-role operators map to
  \texttt{causal.*} procedures over the compiled snapshot (\S\ref{sec:index});
  d-separation is exposed as \texttt{causal.dsep} and is the inner test that
  identification issues repeatedly. Definition statements compile to Cypher
  writes plus a snapshot recompile.}
  \label{tab:operators}
  \small
  \setlength{\tabcolsep}{4pt}
  \begin{tabular}{@{}p{0.30\linewidth} p{0.34\linewidth} p{0.28\linewidth}@{}}
    \toprule
    \textbf{DSL statement} & \textbf{Graphical meaning} & \textbf{Physical operator} \\
    \midrule
    \multicolumn{3}{@{}l}{\textit{Definition}}\\
    \lstinline[style=causaldsl]|DEFINE ... DAG| & write / update model & \texttt{MERGE} + \texttt{causal.compile} \\
    \addlinespace[2pt]
    \multicolumn{3}{@{}l}{\textit{Reachability}}\\
    \lstinline[style=causaldsl]|CAUSES OF X|       & $\mathrm{pa}(X)$ & \texttt{causal.parents} \\
    \lstinline[style=causaldsl]|EFFECTS OF X|      & $\mathrm{ch}(X)$ & \texttt{causal.children} \\
    \lstinline[style=causaldsl]|ANCESTORS OF X|    & $\mathrm{an}(X)$ & \texttt{causal.ancestors} \\
    \lstinline[style=causaldsl]|DESCENDANTS OF X|  & $\mathrm{de}(X)$ & \texttt{causal.descendants} \\
    \addlinespace[2pt]
    \multicolumn{3}{@{}l}{\textit{Path roles}}\\
    \lstinline[style=causaldsl]|MEDIATORS ...|   & on a directed $T\!\to\!Y$ path & \texttt{causal.mediators} \\
    \lstinline[style=causaldsl]|COLLIDERS ...|   & $\to Z \leftarrow$ on a path & \texttt{causal.colliders} \\
    \lstinline[style=causaldsl]|CONFOUNDERS ...| & $\mathrm{an}(T)\cap\mathrm{an}(Y)$ & \texttt{causal.confounders} \\
    \lstinline[style=causaldsl]|BACKDOOR PATHS ...| & path with edge into $T$ & \texttt{causal.backdoorVariables} \\
    \addlinespace[2pt]
    \multicolumn{3}{@{}l}{\textit{Independence \& identification}}\\
    (inner test) & $X\!\perp\!Y\mid Z$ (d-sep) & \texttt{causal.dsep} \\
    \lstinline[style=causaldsl]|IDENTIFY EFFECT ...| & backdoor criterion, adj.\ sets & \texttt{causal.adjustmentCandidate} \\
    \bottomrule
  \end{tabular}
\end{table}

Structural operators map to procedures one-to-one. The reachability operators
compile to \texttt{causal.parents}, \texttt{causal.children},
\texttt{causal.ancestors}, and \texttt{causal.descendants}; the path-role
operators to \texttt{causal.mediators}, \texttt{causal.colliders}, and
\texttt{causal.confounders}; \lstinline[style=causaldsl]|BACKDOOR PATHS| to
\texttt{causal.backdoorVariables}; and \lstinline[style=causaldsl]|IDENTIFY EFFECT|
to \texttt{causal.adjustmentCandidate}, which returns whether an adjustment set
exists, the candidate sets, and---under \lstinline[style=causaldsl]|EXPLAIN|---an
explanation. D-separation is not a top-level keyword but the inner test that
identification issues repeatedly; it is exposed as the \texttt{causal.dsep}
procedure so that the cost of the traversal-intensive query (R2) is borne by a
single indexed operator (\S\ref{sec:index}) rather than by client-side path
enumeration. The quantitative statements (\lstinline[style=causaldsl]|INTERVENE|,
\lstinline[style=causaldsl]|COUNTERFACTUAL|) are handled by the quantitative
layer of \S\ref{sec:quant}; numeric intervention compiles to
\texttt{causal.linearGaussianInference}, while world construction is resolved in
the registry described next.

\subsection{Model registry and causal worlds}\label{sec:lang:registry}

A causal database is rarely a single fixed graph: a curator maintains several
models, and an intervention $do(X{=}x)$ defines a \emph{derived} model---the same
DAG with the edges into $X$ removed---that an analyst wants to name, query, and
compare against the original. The language therefore addresses models by name.
\lstinline[style=causaldsl]|DEFINE MODEL| registers a base model;
\lstinline[style=causaldsl]|USE MODEL| makes a named model (base or derived) the
active graph against which subsequent structural statements run; and
\lstinline[style=causaldsl]|RESET| returns to the base model while leaving derived
worlds registered and addressable. Context switching is resolved entirely in the
registry and emits no Cypher, so moving between worlds costs a name lookup rather
than a re-read of the graph. This registry is also the seam at which the
quantitative layer attaches: an \lstinline[style=causaldsl]|INTERVENE| statement
materializes its derived world once and registers it under a canonical name, so
that a later \lstinline[style=causaldsl]|USE MODEL| over that world reuses the
same compiled snapshot. The mechanics of snapshotting a model---base or
derived---into the layout the procedures read is the subject of the next section.

%% Requires the `causaldsl' listings style and the \TODO/\NOTE macros from main.tex.
%% Table label tab:operators is referenced from sections/background.tex (\S2.2).
%% Operator->procedure mapping is faithful to dsl/.../exec/CypherTranslator.java
%% and plugin/.../api/ProcedureNames.java.
%% ============================================================================
\section{The Causal Query Language}\label{sec:lang}

Requirement R1 asks for a declarative interface in which each catalog query is a
single statement, phrased close to how an analyst states a causal question, that
compiles to physical operators over the graph store. This section presents that
language: its surface (\S\ref{sec:lang:surface}), the mapping from each operator
to the physical operator that answers it (\S\ref{sec:lang:compile},
Table~\ref{tab:operators}), and the model registry that lets a single database
hold several causal worlds at once (\S\ref{sec:lang:registry}). The physical
operators themselves---and why they are fast---are the subject of
\S\ref{sec:index}; here we are concerned only with the interface and its
compilation.

\subsection{Design and surface syntax}\label{sec:lang:surface}

The language is a small statement-oriented DSL. A program is a sequence of
statements; each statement is one of three kinds. \emph{Definition} statements
declare variables and edges and thereby write the causal model into the graph.
\emph{Structural} statements ask the catalog queries of \S\ref{sec:bg:catalog}:
reachability, path roles, and identification. \emph{Quantitative} statements
express interventions and counterfactuals over a parameterized model
(\S\ref{sec:quant}). Two design choices follow directly from R1. First, an
operator names a causal concept, not a traversal: an analyst writes
\lstinline[style=causaldsl]|CONFOUNDERS OF T AND Y| rather than a path pattern,
so the statement reads as the question and its compilation to a graph operation
is the system's concern, not the user's. Second, every statement is
self-contained and refers to variables by their canonical
\textsf{id}, so a query is portable across any database holding a model with
those variables.

Definition statements come in a standalone and a block form. A bare
\lstinline[style=causaldsl]|DEFINE CAUSAL VARIABLE| introduces a single variable
with its observation status; a \lstinline[style=causaldsl]|DEFINE CAUSAL DAG|
block declares a model as a set of edges, optionally carrying a structural
coefficient per edge for the quantitative layer:

\begin{lstlisting}[style=causaldsl]
DEFINE CAUSAL DAG CovidModel {
  covid19_burden    -> interventions
  covid19_burden    -> exposure_sarscov2
  interventions     -> mobility
  mobility          -> exposure_sarscov2
  exposure_sarscov2 -> onset_covid19
  onset_covid19     -> reported_cases
}
\end{lstlisting}

\noindent
Against the model so defined, the structural questions of our running example
(\S\ref{sec:bg:example}) are each one statement:

\begin{lstlisting}[style=causaldsl]
CONFOUNDERS OF interventions AND reported_cases
BACKDOOR PATHS BETWEEN interventions AND reported_cases
IDENTIFY EFFECT OF interventions ON reported_cases
  RETURN ADJUSTMENT SETS EXPLAIN
\end{lstlisting}

\noindent
The first returns \textsf{covid19\_burden}; the second reports the backdoor path
it opens; the third decides identifiability and, with
\lstinline[style=causaldsl]|RETURN ADJUSTMENT SETS|, returns the valid adjustment
sets, the optional \lstinline[style=causaldsl]|EXPLAIN| attaching a
human-readable justification. The reachability and path-role operators
(\lstinline[style=causaldsl]|CAUSES OF|, \lstinline[style=causaldsl]|EFFECTS OF|,
\lstinline[style=causaldsl]|ANCESTORS OF|, \lstinline[style=causaldsl]|DESCENDANTS OF|,
\lstinline[style=causaldsl]|MEDIATORS BETWEEN|, \lstinline[style=causaldsl]|COLLIDERS BETWEEN|)
follow the same one-statement-per-question shape. The full grammar is given in
ANTLR form in the artifact; Table~\ref{tab:operators} lists the operators that
matter for what follows.

\subsection{Compilation to physical operators}\label{sec:lang:compile}

A statement is parsed to an AST and then compiled, statement by statement, to one
or more Cypher units that execute inside the DBMS. The compiler emits exactly two
kinds of unit. A \emph{definition} statement compiles to Cypher \emph{writes}
that maintain the schema of \S\ref{sec:bg:dag}: a \lstinline[style=causaldsl]|DEFINE|
becomes a \texttt{MERGE} of the \textsf{Variable} nodes and \textsf{CAUSES} edges
it introduces (a single batched \texttt{UNWIND} for the edge set), followed by a
call to \texttt{causal.compile} that refreshes the compiled snapshot
(\S\ref{sec:compile}). A \emph{structural} statement compiles to a read: a
\texttt{CALL} of the \texttt{causal.*} procedure that answers it, which consults
the snapshot rather than re-traversing the stored graph. This two-target
compilation is what keeps the model \emph{in} the database: the language never
holds the graph in client memory, and the same authoritative snapshot serves
every reader.

\begin{table}[t]
  \centering
  \caption{Language operators, their graphical semantics, and the physical
  operator each compiles to. Reachability and path-role operators map to
  \texttt{causal.*} procedures over the compiled snapshot (\S\ref{sec:index});
  d-separation is exposed as \texttt{causal.dsep} and is the inner test that
  identification issues repeatedly. Definition statements compile to Cypher
  writes plus a snapshot recompile.}
  \label{tab:operators}
  \small
  \setlength{\tabcolsep}{4pt}
  \begin{tabular}{@{}p{0.30\linewidth} p{0.34\linewidth} p{0.28\linewidth}@{}}
    \toprule
    \textbf{DSL statement} & \textbf{Graphical meaning} & \textbf{Physical operator} \\
    \midrule
    \multicolumn{3}{@{}l}{\textit{Definition}}\\
    \lstinline[style=causaldsl]|DEFINE ... DAG| & write / update model & \texttt{MERGE} + \texttt{causal.compile} \\
    \addlinespace[2pt]
    \multicolumn{3}{@{}l}{\textit{Reachability}}\\
    \lstinline[style=causaldsl]|CAUSES OF X|       & $\mathrm{pa}(X)$ & \texttt{causal.parents} \\
    \lstinline[style=causaldsl]|EFFECTS OF X|      & $\mathrm{ch}(X)$ & \texttt{causal.children} \\
    \lstinline[style=causaldsl]|ANCESTORS OF X|    & $\mathrm{an}(X)$ & \texttt{causal.ancestors} \\
    \lstinline[style=causaldsl]|DESCENDANTS OF X|  & $\mathrm{de}(X)$ & \texttt{causal.descendants} \\
    \addlinespace[2pt]
    \multicolumn{3}{@{}l}{\textit{Path roles}}\\
    \lstinline[style=causaldsl]|MEDIATORS ...|   & on a directed $T\!\to\!Y$ path & \texttt{causal.mediators} \\
    \lstinline[style=causaldsl]|COLLIDERS ...|   & $\to Z \leftarrow$ on a path & \texttt{causal.colliders} \\
    \lstinline[style=causaldsl]|CONFOUNDERS ...| & $\mathrm{an}(T)\cap\mathrm{an}(Y)$ & \texttt{causal.confounders} \\
    \lstinline[style=causaldsl]|BACKDOOR PATHS ...| & path with edge into $T$ & \texttt{causal.backdoorVariables} \\
    \addlinespace[2pt]
    \multicolumn{3}{@{}l}{\textit{Independence \& identification}}\\
    (inner test) & $X\!\perp\!Y\mid Z$ (d-sep) & \texttt{causal.dsep} \\
    \lstinline[style=causaldsl]|IDENTIFY EFFECT ...| & backdoor criterion, adj.\ sets & \texttt{causal.adjustmentCandidate} \\
    \bottomrule
  \end{tabular}
\end{table}

Structural operators map to procedures one-to-one. The reachability operators
compile to \texttt{causal.parents}, \texttt{causal.children},
\texttt{causal.ancestors}, and \texttt{causal.descendants}; the path-role
operators to \texttt{causal.mediators}, \texttt{causal.colliders}, and
\texttt{causal.confounders}; \lstinline[style=causaldsl]|BACKDOOR PATHS| to
\texttt{causal.backdoorVariables}; and \lstinline[style=causaldsl]|IDENTIFY EFFECT|
to \texttt{causal.adjustmentCandidate}, which returns whether an adjustment set
exists, the candidate sets, and---under \lstinline[style=causaldsl]|EXPLAIN|---an
explanation. D-separation is not a top-level keyword but the inner test that
identification issues repeatedly; it is exposed as the \texttt{causal.dsep}
procedure so that the cost of the traversal-intensive query (R2) is borne by a
single indexed operator (\S\ref{sec:index}) rather than by client-side path
enumeration. The quantitative statements (\lstinline[style=causaldsl]|INTERVENE|,
\lstinline[style=causaldsl]|COUNTERFACTUAL|) are handled by the quantitative
layer of \S\ref{sec:quant}; numeric intervention compiles to
\texttt{causal.linearGaussianInference}, while world construction is resolved in
the registry described next.

\subsection{Model registry and causal worlds}\label{sec:lang:registry}

A causal database is rarely a single fixed graph: a curator maintains several
models, and an intervention $do(X{=}x)$ defines a \emph{derived} model---the same
DAG with the edges into $X$ removed---that an analyst wants to name, query, and
compare against the original. The language therefore addresses models by name.
\lstinline[style=causaldsl]|DEFINE MODEL| registers a base model;
\lstinline[style=causaldsl]|USE MODEL| makes a named model (base or derived) the
active graph against which subsequent structural statements run; and
\lstinline[style=causaldsl]|RESET| returns to the base model while leaving derived
worlds registered and addressable. Context switching is resolved entirely in the
registry and emits no Cypher, so moving between worlds costs a name lookup rather
than a re-read of the graph. This registry is also the seam at which the
quantitative layer attaches: an \lstinline[style=causaldsl]|INTERVENE| statement
materializes its derived world once and registers it under a canonical name, so
that a later \lstinline[style=causaldsl]|USE MODEL| over that world reuses the
same compiled snapshot. The mechanics of snapshotting a model---base or
derived---into the layout the procedures read is the subject of the next section.

%% Requires the `causaldsl' listings style and the \TODO/\NOTE macros from main.tex.
%% Table label tab:operators is referenced from sections/background.tex (\S2.2).
%% Operator->procedure mapping is faithful to dsl/.../exec/CypherTranslator.java
%% and plugin/.../api/ProcedureNames.java.
%% ============================================================================
\section{The Causal Query Language}\label{sec:lang}

Requirement R1 asks for a declarative interface in which each catalog query is a
single statement, phrased close to how an analyst states a causal question, that
compiles to physical operators over the graph store. This section presents that
language: its surface (\S\ref{sec:lang:surface}), the mapping from each operator
to the physical operator that answers it (\S\ref{sec:lang:compile},
Table~\ref{tab:operators}), and the model registry that lets a single database
hold several causal worlds at once (\S\ref{sec:lang:registry}). The physical
operators themselves---and why they are fast---are the subject of
\S\ref{sec:index}; here we are concerned only with the interface and its
compilation.

\subsection{Design and surface syntax}\label{sec:lang:surface}

The language is a small statement-oriented DSL. A program is a sequence of
statements; each statement is one of three kinds. \emph{Definition} statements
declare variables and edges and thereby write the causal model into the graph.
\emph{Structural} statements ask the catalog queries of \S\ref{sec:bg:catalog}:
reachability, path roles, and identification. \emph{Quantitative} statements
express interventions and counterfactuals over a parameterized model
(\S\ref{sec:quant}). Two design choices follow directly from R1. First, an
operator names a causal concept, not a traversal: an analyst writes
\lstinline[style=causaldsl]|CONFOUNDERS OF T AND Y| rather than a path pattern,
so the statement reads as the question and its compilation to a graph operation
is the system's concern, not the user's. Second, every statement is
self-contained and refers to variables by their canonical
\textsf{id}, so a query is portable across any database holding a model with
those variables.

Definition statements come in a standalone and a block form. A bare
\lstinline[style=causaldsl]|DEFINE CAUSAL VARIABLE| introduces a single variable
with its observation status; a \lstinline[style=causaldsl]|DEFINE CAUSAL DAG|
block declares a model as a set of edges, optionally carrying a structural
coefficient per edge for the quantitative layer:

\begin{lstlisting}[style=causaldsl]
DEFINE CAUSAL DAG CovidModel {
  covid19_burden    -> interventions
  covid19_burden    -> exposure_sarscov2
  interventions     -> mobility
  mobility          -> exposure_sarscov2
  exposure_sarscov2 -> onset_covid19
  onset_covid19     -> reported_cases
}
\end{lstlisting}

\noindent
Against the model so defined, the structural questions of our running example
(\S\ref{sec:bg:example}) are each one statement:

\begin{lstlisting}[style=causaldsl]
CONFOUNDERS OF interventions AND reported_cases
BACKDOOR PATHS BETWEEN interventions AND reported_cases
IDENTIFY EFFECT OF interventions ON reported_cases
  RETURN ADJUSTMENT SETS EXPLAIN
\end{lstlisting}

\noindent
The first returns \textsf{covid19\_burden}; the second reports the backdoor path
it opens; the third decides identifiability and, with
\lstinline[style=causaldsl]|RETURN ADJUSTMENT SETS|, returns the valid adjustment
sets, the optional \lstinline[style=causaldsl]|EXPLAIN| attaching a
human-readable justification. The reachability and path-role operators
(\lstinline[style=causaldsl]|CAUSES OF|, \lstinline[style=causaldsl]|EFFECTS OF|,
\lstinline[style=causaldsl]|ANCESTORS OF|, \lstinline[style=causaldsl]|DESCENDANTS OF|,
\lstinline[style=causaldsl]|MEDIATORS BETWEEN|, \lstinline[style=causaldsl]|COLLIDERS BETWEEN|)
follow the same one-statement-per-question shape. The full grammar is given in
ANTLR form in the artifact; Table~\ref{tab:operators} lists the operators that
matter for what follows.

\subsection{Compilation to physical operators}\label{sec:lang:compile}

A statement is parsed to an AST and then compiled, statement by statement, to one
or more Cypher units that execute inside the DBMS. The compiler emits exactly two
kinds of unit. A \emph{definition} statement compiles to Cypher \emph{writes}
that maintain the schema of \S\ref{sec:bg:dag}: a \lstinline[style=causaldsl]|DEFINE|
becomes a \texttt{MERGE} of the \textsf{Variable} nodes and \textsf{CAUSES} edges
it introduces (a single batched \texttt{UNWIND} for the edge set), followed by a
call to \texttt{causal.compile} that refreshes the compiled snapshot
(\S\ref{sec:compile}). A \emph{structural} statement compiles to a read: a
\texttt{CALL} of the \texttt{causal.*} procedure that answers it, which consults
the snapshot rather than re-traversing the stored graph. This two-target
compilation is what keeps the model \emph{in} the database: the language never
holds the graph in client memory, and the same authoritative snapshot serves
every reader.

\begin{table}[t]
  \centering
  \caption{Language operators, their graphical semantics, and the physical
  operator each compiles to. Reachability and path-role operators map to
  \texttt{causal.*} procedures over the compiled snapshot (\S\ref{sec:index});
  d-separation is exposed as \texttt{causal.dsep} and is the inner test that
  identification issues repeatedly. Definition statements compile to Cypher
  writes plus a snapshot recompile.}
  \label{tab:operators}
  \small
  \setlength{\tabcolsep}{4pt}
  \begin{tabular}{@{}p{0.30\linewidth} p{0.34\linewidth} p{0.28\linewidth}@{}}
    \toprule
    \textbf{DSL statement} & \textbf{Graphical meaning} & \textbf{Physical operator} \\
    \midrule
    \multicolumn{3}{@{}l}{\textit{Definition}}\\
    \lstinline[style=causaldsl]|DEFINE ... DAG| & write / update model & \texttt{MERGE} + \texttt{causal.compile} \\
    \addlinespace[2pt]
    \multicolumn{3}{@{}l}{\textit{Reachability}}\\
    \lstinline[style=causaldsl]|CAUSES OF X|       & $\mathrm{pa}(X)$ & \texttt{causal.parents} \\
    \lstinline[style=causaldsl]|EFFECTS OF X|      & $\mathrm{ch}(X)$ & \texttt{causal.children} \\
    \lstinline[style=causaldsl]|ANCESTORS OF X|    & $\mathrm{an}(X)$ & \texttt{causal.ancestors} \\
    \lstinline[style=causaldsl]|DESCENDANTS OF X|  & $\mathrm{de}(X)$ & \texttt{causal.descendants} \\
    \addlinespace[2pt]
    \multicolumn{3}{@{}l}{\textit{Path roles}}\\
    \lstinline[style=causaldsl]|MEDIATORS ...|   & on a directed $T\!\to\!Y$ path & \texttt{causal.mediators} \\
    \lstinline[style=causaldsl]|COLLIDERS ...|   & $\to Z \leftarrow$ on a path & \texttt{causal.colliders} \\
    \lstinline[style=causaldsl]|CONFOUNDERS ...| & $\mathrm{an}(T)\cap\mathrm{an}(Y)$ & \texttt{causal.confounders} \\
    \lstinline[style=causaldsl]|BACKDOOR PATHS ...| & path with edge into $T$ & \texttt{causal.backdoorVariables} \\
    \addlinespace[2pt]
    \multicolumn{3}{@{}l}{\textit{Independence \& identification}}\\
    (inner test) & $X\!\perp\!Y\mid Z$ (d-sep) & \texttt{causal.dsep} \\
    \lstinline[style=causaldsl]|IDENTIFY EFFECT ...| & backdoor criterion, adj.\ sets & \texttt{causal.adjustmentCandidate} \\
    \bottomrule
  \end{tabular}
\end{table}

Structural operators map to procedures one-to-one. The reachability operators
compile to \texttt{causal.parents}, \texttt{causal.children},
\texttt{causal.ancestors}, and \texttt{causal.descendants}; the path-role
operators to \texttt{causal.mediators}, \texttt{causal.colliders}, and
\texttt{causal.confounders}; \lstinline[style=causaldsl]|BACKDOOR PATHS| to
\texttt{causal.backdoorVariables}; and \lstinline[style=causaldsl]|IDENTIFY EFFECT|
to \texttt{causal.adjustmentCandidate}, which returns whether an adjustment set
exists, the candidate sets, and---under \lstinline[style=causaldsl]|EXPLAIN|---an
explanation. D-separation is not a top-level keyword but the inner test that
identification issues repeatedly; it is exposed as the \texttt{causal.dsep}
procedure so that the cost of the traversal-intensive query (R2) is borne by a
single indexed operator (\S\ref{sec:index}) rather than by client-side path
enumeration. The quantitative statements (\lstinline[style=causaldsl]|INTERVENE|,
\lstinline[style=causaldsl]|COUNTERFACTUAL|) are handled by the quantitative
layer of \S\ref{sec:quant}; numeric intervention compiles to
\texttt{causal.linearGaussianInference}, while world construction is resolved in
the registry described next.

\subsection{Model registry and causal worlds}\label{sec:lang:registry}

A causal database is rarely a single fixed graph: a curator maintains several
models, and an intervention $do(X{=}x)$ defines a \emph{derived} model---the same
DAG with the edges into $X$ removed---that an analyst wants to name, query, and
compare against the original. The language therefore addresses models by name.
\lstinline[style=causaldsl]|DEFINE MODEL| registers a base model;
\lstinline[style=causaldsl]|USE MODEL| makes a named model (base or derived) the
active graph against which subsequent structural statements run; and
\lstinline[style=causaldsl]|RESET| returns to the base model while leaving derived
worlds registered and addressable. Context switching is resolved entirely in the
registry and emits no Cypher, so moving between worlds costs a name lookup rather
than a re-read of the graph. This registry is also the seam at which the
quantitative layer attaches: an \lstinline[style=causaldsl]|INTERVENE| statement
materializes its derived world once and registers it under a canonical name, so
that a later \lstinline[style=causaldsl]|USE MODEL| over that world reuses the
same compiled snapshot. The mechanics of snapshotting a model---base or
derived---into the layout the procedures read is the subject of the next section.
